\begin{helper}
Fix \(n,m,k\), an edge density \(p\), and parameters
\(\varepsilon,\varepsilon_1,\varepsilon_2\).  Let \(\iota\) be a finite
candidate index type, and let \(C_i(\omega)\) be candidate events.  Suppose
that the bad-on-regular event
\[
\neg\noderef{TwoBiteMediumClosedPairsEvent}_{k,\varepsilon,\varepsilon_1}
  (\omega)
\quad\text{and}\quad
\noderef{FiberAndDegreeEvent}(\omega,\varepsilon_2)
\]
is covered by the candidates: every sample satisfying it lies in some
\(C_i\).  Suppose also that, after conditioning on each embedding cell
\(\noderef{TwoBiteEmbedding}(\omega)=\pi\), each candidate satisfies
\[
\noderef{TwoBiteEventProbability}
  (n,m,p,\{\omega:\noderef{TwoBiteEmbedding}(\omega)=\pi\ \text{and}\ C_i(\omega)\})
\le
e^{-T}\,
\noderef{TwoBiteEventProbability}
  (n,m,p,\{\omega:\noderef{TwoBiteEmbedding}(\omega)=\pi\})
\]
for a tail exponent \(T\).  If the number of candidates is at most \(e^C\),
then
\[
\noderef{TwoBiteEventProbability}
  (n,m,p,\{\omega:
    \neg\noderef{TwoBiteMediumClosedPairsEvent}_{k,\varepsilon,\varepsilon_1}
      (\omega)
    \ \text{and}\
    \noderef{FiberAndDegreeEvent}(\omega,\varepsilon_2)\})
\le e^{C-T}.
\]
\end{helper}
\begin{proof}
For a fixed candidate \(i\), apply
\(\noderef{TwoBiteEventProbabilitySumOverEmbeddings}\) to the event \(C_i\).
The displayed embedding-cell estimate is exactly the required per-cell
inequality with the bound \(e^{-T}\), and \(e^{-T}\) is nonnegative.  Hence
\[
\noderef{TwoBiteEventProbability}(n,m,p,C_i)\le e^{-T}
\]
for every \(i\).

Now apply \(\noderef{MediumClosedPairsFiniteCandidateUnionBound}\) to the
bad-on-regular event and the finite family \((C_i)_{i\in\iota}\).  The cover
hypothesis gives the event inclusion needed by that lemma, and \(0\le p\le1\)
gives nonnegativity of the sample weights.  Therefore
\[
\noderef{TwoBiteEventProbability}(\text{bad-on-regular})
\le |\iota|\,e^{-T}.
\]
Using \(|\iota|\le e^C\) and \(e^{-T}\ge0\), this is at most
\[
e^C e^{-T}=e^{C-T},
\]
which proves the claimed finite-\(n\) exponential upper bound.
\end{proof}
